% GENERATED FILE. Edit canonical lessons, then run npm run generate:books.
% canonical-source-hash: fnv1a64:beb3127d5a197453
% canonical-lessons: BN-C252-kakima, BN-C252-bhaipo, BN-C252-nati, BN-C252-natni, BN-C252-shvoshur

\chapter{Paternal aunt, Nephew, Grandson, Granddaughter, Father-in-law}
\label{ch:bn-paternal-aunt-nephew-grandson-granddaughter-father-in-law}
\hlchaptermodality{252}

\begin{chapteropening}
\textbf{By the end of this chapter:} \emph{I can say a paternal aunt, a nephew, a grandson, a granddaughter and a father-in-law.}

Complete the last lesson of chapter 252: I can say a paternal aunt, a nephew, a grandson, a granddaughter and a father-in-law.
\end{chapteropening}

\section[\texorpdfstring{\bn{কাকিমা} (kākimā)}{কাকিমা (kākimā)}]{\bn{কাকিমা} (kākimā) — a paternal aunt (by marriage)}

\label{lesson:BN-C252-kakima}



\begin{quote}
Before the new one: say the Bengali for a party, then the Bengali for a journey.
\end{quote}

\subsection*{You'll want to know: \bn{কাকিমা}}
\textbf{\bn{কাকিমা}} — \emph{kākimā} — \textquotedblleft{}a paternal aunt (by marriage)\textquotedblright{}.

\begin{grammarlens}[title={What you've built}]
One more word for this chapter.
\end{grammarlens}

\subsection*{Your turn}
\begin{itemize}
\raggedright
  \item \emph{Say it:} \emph{kākimā}
  \item \emph{Say it:} \emph{kākimā}, once more
  \item \emph{Say it:} say \emph{kākimā}
  \item \emph{Recall:} say the Bengali for a party, then the Bengali for a journey, then say \emph{kākimā} again
\end{itemize}

\begin{tcolorbox}[breakable,colback=teal!4,colframe=teal!35!black,title={Before you move on}]
What does \bn{কাকিমা} mean? (\textquotedblleft{}A paternal aunt (by marriage)\textquotedblright{}.) Say it once more.
\end{tcolorbox}
\section[\texorpdfstring{\bn{ভাইপো} (bhāipo)}{ভাইপো (bhāipo)}]{\bn{ভাইপো} (bhāipo) — a nephew (brother's son)}

\label{lesson:BN-C252-bhaipo}



\begin{quote}
Before the new one: say the Bengali for a journey, then the Bengali for a paternal aunt.
\end{quote}

\subsection*{You'll want to know: \bn{ভাইপো}}
\textbf{\bn{ভাইপো}} — \emph{bhāipo} — \textquotedblleft{}a nephew (brother's son)\textquotedblright{}.

\begin{grammarlens}[title={What you've built}]
One more word for this chapter.
\end{grammarlens}

\subsection*{Your turn}
\begin{itemize}
\raggedright
  \item \emph{Say it:} \emph{bhāipo}
  \item \emph{Say it:} \emph{bhāipo}, once more
  \item \emph{Say it:} say \emph{bhāipo}
  \item \emph{Recall:} say the Bengali for a journey, then the Bengali for a paternal aunt, then say \emph{bhāipo} again
\end{itemize}

\begin{tcolorbox}[breakable,colback=teal!4,colframe=teal!35!black,title={Before you move on}]
What does \bn{ভাইপো} mean? (\textquotedblleft{}A nephew (brother's son)\textquotedblright{}.) Say it once more.
\end{tcolorbox}
\section[\texorpdfstring{\bn{নাতি} (nāti)}{নাতি (nāti)}]{\bn{নাতি} (nāti) — a grandson}

\label{lesson:BN-C252-nati}



\begin{quote}
Before the new one: say the Bengali for a paternal aunt, then the Bengali for a nephew.
\end{quote}

\subsection*{You'll want to know: \bn{নাতি}}
\textbf{\bn{নাতি}} — \emph{nāti} — \textquotedblleft{}a grandson\textquotedblright{}.

\begin{grammarlens}[title={What you've built}]
One more word for this chapter.
\end{grammarlens}

\subsection*{Your turn}
\begin{itemize}
\raggedright
  \item \emph{Say it:} \emph{nāti}
  \item \emph{Say it:} \emph{nāti}, once more
  \item \emph{Say it:} say \emph{nāti}
  \item \emph{Recall:} say the Bengali for a paternal aunt, then the Bengali for a nephew, then say \emph{nāti} again
\end{itemize}

\begin{tcolorbox}[breakable,colback=teal!4,colframe=teal!35!black,title={Before you move on}]
What does \bn{নাতি} mean? (\textquotedblleft{}A grandson\textquotedblright{}.) Say it once more.
\end{tcolorbox}
\section[\texorpdfstring{\bn{নাতনি} (nātni)}{নাতনি (nātni)}]{\bn{নাতনি} (nātni) — a granddaughter}

\label{lesson:BN-C252-natni}



\begin{quote}
Before the new one: say the Bengali for a nephew, then the Bengali for a grandson.
\end{quote}

\subsection*{You'll want to know: \bn{নাতনি}}
\textbf{\bn{নাতনি}} — \emph{nātni} — \textquotedblleft{}a granddaughter\textquotedblright{}.

\begin{grammarlens}[title={What you've built}]
One more word for this chapter.
\end{grammarlens}

\subsection*{Your turn}
\begin{itemize}
\raggedright
  \item \emph{Say it:} \emph{nātni}
  \item \emph{Say it:} \emph{nātni}, once more
  \item \emph{Say it:} say \emph{nātni}
  \item \emph{Recall:} say the Bengali for a nephew, then the Bengali for a grandson, then say \emph{nātni} again
\end{itemize}

\begin{tcolorbox}[breakable,colback=teal!4,colframe=teal!35!black,title={Before you move on}]
What does \bn{নাতনি} mean? (\textquotedblleft{}A granddaughter\textquotedblright{}.) Say it once more.
\end{tcolorbox}
\section[\texorpdfstring{\bn{শ্বশুর} (shvôshur)}{শ্বশুর (shvôshur)}]{\bn{শ্বশুর} (shvôshur) — a father-in-law}

\label{lesson:BN-C252-shvoshur}



\begin{quote}
Before the new one: say the Bengali for a grandson, then the Bengali for a granddaughter.
\end{quote}

\subsection*{You'll want to know: \bn{শ্বশুর}}
\textbf{\bn{শ্বশুর}} — \emph{shvôshur} — \textquotedblleft{}a father-in-law\textquotedblright{}.

\begin{grammarlens}[title={What you've built}]
One more word for this chapter.
\end{grammarlens}

\subsection*{Your turn}
\begin{itemize}
\raggedright
  \item \emph{Say it:} \emph{shvôshur}
  \item \emph{Say it:} \emph{shvôshur}, once more
  \item \emph{Say it:} say \emph{shvôshur}
  \item \emph{Recall:} say the Bengali for a grandson, then the Bengali for a granddaughter, then say \emph{shvôshur} again
\end{itemize}

\begin{tcolorbox}[breakable,colback=teal!4,colframe=teal!35!black,title={Before you move on}]
What does \bn{শ্বশুর} mean? (\textquotedblleft{}A father-in-law\textquotedblright{}.) Say it once more.
\end{tcolorbox}
